\title{符号圈平方图中的半胞通量与可精确求解的八周期相}
\author{Yicheng Zhao \and Jiachen Li}
\date{}
\maketitle

\begin{center}
\small
重庆邮电大学数学与统计学院，中国重庆\\[2pt]
Yicheng Zhao：\texttt{2023213805@cqupt.edu.cn}，
ORCID \href{https://orcid.org/0009-0003-7618-1661}{0009-0003-7618-1661}\\
Jiachen Li：\texttt{2023213809@stu.cqupt.edu.cn}，
ORCID \href{https://orcid.org/0009-0006-5119-3369}{0009-0006-5119-3369}
\end{center}

\begin{abstract}
在固定底图上改变边的符号，不会改变邻接矩阵的支撑，却会改变其谱。本文研究圈平方图
$C_n(1,2)$ 上的这一谱优化问题。对于任意偶周期 Hamilton 规范符号，我们证明一个
switching 不变的半胞通量条件等价于自然的单项式 chiral 对合，从而把 $2m$ 维
Floquet 问题约化为 $m$ 维平方谱问题。平方 Bloch 谱边小于 $8$ 的最早实现出现在
本原周期八。对相应的对径双缺陷相，我们求出全部四条平方色散支，并证明对每个
$L\geq1$，
\[
 \rho(A_{8L,+})^2=4+\sqrt{10+2\sqrt5}.
\]
负 holonomy 有限环也具有显式谱半径。在本原周期八内，该相在自然对称意义下是唯一的
sub-eight 相，而所有更短本原周期都不能达到 sub-eight。特别地，对每个 $L\geq4$，
这一正 holonomy 相的谱半径严格小于 $C_{8L}(1,2)$ 上自然的 twisted 符号。
\end{abstract}

\noindent\textbf{关键词：} 符号图；圈平方图；谱半径；switching；Floquet 分解；chiral 对称

\medskip
\noindent\textbf{MSC 2020：} 05C50（主要）；05C22，15A18（次要）
